\documentclass[12pt]{article}
\usepackage{amsmath,amssymb,amsthm}
\usepackage[margin=1in]{geometry}
\usepackage{algorithm}
\usepackage{algpseudocode}

\newtheorem{theorem}{Theorem}
\newtheorem{lemma}[theorem]{Lemma}

\title{Problem 201: Subsets with a Unique Sum}
\author{Project Euler Solution}
\date{}

\begin{document}
\maketitle

\section{Problem Statement}
For any set $A$ of numbers, let $\operatorname{sum}(A)$ denote the sum of elements of $A$.
Consider the set $S = \{1^2, 2^2, \ldots, 100^2\}$. Let $U$ be the set of all integers $v$
such that exactly one 50-element subset of $S$ has sum equal to $v$. Find $\operatorname{sum}(U)$.

\section{Mathematical Foundation}

\begin{theorem}[Sum Bounds]
Let $S_{\min}$ and $S_{\max}$ denote the minimum and maximum possible sums of a 50-element
subset of $S$. Then
\[
S_{\min} = \sum_{k=1}^{50} k^2 = \frac{50 \cdot 51 \cdot 101}{6} = 42925, \qquad
S_{\max} = \sum_{k=51}^{100} k^2 = 295425.
\]
\end{theorem}

\begin{proof}
The minimum sum is achieved by choosing the 50 smallest elements $\{1^2, \ldots, 50^2\}$,
and the maximum by choosing the 50 largest $\{51^2, \ldots, 100^2\}$. By the formula
$\sum_{k=1}^{n} k^2 = n(n+1)(2n+1)/6$, we obtain $S_{\min} = 42925$. Then
$S_{\max} = \sum_{k=1}^{100} k^2 - S_{\min} = 338350 - 42925 = 295425$.
\end{proof}

\begin{lemma}[Three-State Sufficiency]
To determine whether a sum value $v$ is achieved by exactly one 50-element subset, it
suffices to track for each $(j, s)$ pair whether the number of $j$-element subsets of the
first $i$ elements summing to $s$ is $0$, $1$, or $\geq 2$.
\end{lemma}

\begin{proof}
We seek values $v$ where the count $c(50, v)$ of 50-element subsets summing to $v$ equals
exactly~1. Since we only need to distinguish among $c = 0$, $c = 1$, and $c \geq 2$,
tracking counts clamped at~2 preserves all necessary information. The recurrence
$c_i(j,s) = c_{i-1}(j,s) + c_{i-1}(j-1, s - a_i)$ is additive; clamping at~2 after each
update is correct because $\min(x+y, 2) = x+y$ when $x+y \leq 1$, and
$\min(x+y, 2) = 2$ when $x+y \geq 2$. Thus the three-state DP correctly identifies all
values with count exactly~1.
\end{proof}

\section{Algorithm}

\begin{algorithmic}[1]
\State Compute $S_{\max} = 295425$
\State Initialize $\mathrm{dp}[0 \ldots k][0 \ldots S_{\max}] \gets 0$; set $\mathrm{dp}[0][0] \gets 1$
\For{$i = 1$ to $n$}
    \State $a_i \gets i^2$
    \For{$j = \min(i, k)$ \textbf{down to} $1$}
        \For{$s = S_{\max}$ \textbf{down to} $a_i$}
            \State $\mathrm{dp}[j][s] \gets \min(\mathrm{dp}[j][s] + \mathrm{dp}[j-1][s - a_i],\; 2)$
        \EndFor
    \EndFor
\EndFor
\State $\mathrm{result} \gets \sum_{s:\;\mathrm{dp}[k][s]=1} s$
\State \Return result
\end{algorithmic}

\section{Complexity Analysis}
\begin{itemize}
    \item \textbf{Time:} $O(n \cdot k \cdot S_{\max}) = O(100 \cdot 50 \cdot 295425) \approx 1.48 \times 10^9$.
    \item \textbf{Space:} $O(k \cdot S_{\max}) = O(50 \cdot 295425) \approx 1.48 \times 10^7$ entries, each storing a value in $\{0,1,2\}$.
\end{itemize}

\section{Answer}

\[
\boxed{115039000}
\]

\end{document}
